% SPDX-License-Identifier: CC-BY-4.0
% Copyright (c) 2026 Martin Schröder <info@swedishembedded.com>

\section{Several generations of a 27B model on its own experience}\label{sec:generations}

The experiments above train a smaller model on a larger one's trajectories, or a small model on
tasks it writes itself. This section asks the question the loop exists for at full scale: does
Qwen3.8-27B, trained only on its own verified attempts, become a better agent generation after
generation, on tasks it never trained on? No stronger model takes part, and the base is the only
starting point.

\subsection{Workloads}

Four kinds of task are attempted by the same model in every generation, each with a training
partition the generation practises on and a held-out partition it is examined on; the partitions
are whole families and are fixed before the base attempts anything.
\begin{itemize}
\item \textbf{Repository faults} (\cref{sec:protocol}): one seeded fault per task in the eight
  authored repositories; the exam is the reserved files of six repositories and every file of the
  two sealed ones (mdrender, graphkit), so the exam measures transfer to repositories never trained on.
\item \textbf{Several faults at once}: \code{splinter tasks combine} joins three admitted
  single-fault tasks of one snapshot whose faults lie in different files into one task (the union
  of their checks and reference changes, one merged report); it is admitted only if no member's
  reference alone repairs it, and a combined task with any reserved file is an exam task.
\item \textbf{Embedded firmware}: defects in small firmware services for Cortex-M and nRF5340
  boards, built with the vendor toolchain and graded by running the image in an emulator against
  hidden behavioural checks.
\item \textbf{Mechanical design}: a part drawing to a manufacturable solid, graded by geometry
  against a reference.
\end{itemize}

\subsection{One generation}

Each generation starts from the champion (the base, then the last promoted adapter) and runs four
steps; every step records its result before the next starts.
\begin{enumerate}
\item \textbf{Explore.} The champion attempts every training task $k$ times ($k=4$ for the
  repository workloads) under the model family's recommended sampling; a separate verifier grades
  each attempt and an exposure audit fails any attempt that reached hidden material.
\item \textbf{Learn.} The graded attempts are admitted as experience (held-out families are
  refused by the exam seal) and become on-policy training records, one per model call. Each
  candidate continues from the champion:
  \begin{itemize}
  \item \emph{GRPO}: an attempt's advantage is its reward (1 verified, 0 not) less its task
    group's mean over the group's standard deviation; the clipped surrogate with a $k_3$ KL to the
    base trains the adapter (clip 0.2, $\beta=0.02$).
  \item \emph{GRPO with privileged self-distillation}: the same, and every task whose attempts all
    failed (a group with no reward signal) is distilled instead: the same base model is shown what
    the verifier knows (the failing checks and their output, the reference change where the task
    has one) and each token the student wrote gets the advantage
    $A_t=\mathrm{clip}(\log q(y_t\mid\text{privileged},y_{<t})-\log\pi_{\text{old}}(y_t\mid y_{<t}),\pm 5)$,
    with a cap of 40 on its sum over a sequence, through the same clipped step.
  \end{itemize}
\item \textbf{Examine.} Every candidate and the champion attempt every held-out task under the
  same agent build, serving build, benchmark revision and load per card.
\item \textbf{Decide.} Each candidate is paired task by task with the champion on every workload;
  the candidate with the largest gain over all paired tasks replaces the champion when that gain is
  positive and no workload's mean falls by more than 0.2. The decision, the paired counts and the
  sign test are recorded whether or not the champion changes.
\end{enumerate}

\subsection{What the harness must not decide}\label{sec:generations-harness}

A comparison of agents is only a comparison of models if the software around them fails neither
arm. Two failures of that kind were found while running the loop and are part of the protocol now.
First, a wall-clock limit per trial makes the server's load part of the outcome: with about a
hundred agents on one card, per-stream decoding slowed until correct trajectories were cut off, and
the base's held-out repository pass rate read far below its rate on the same exam at a fifth of the
load. The loop therefore gives a model examined on a share of the cards that share of a workload's trials.
Second, a serving engine that ran out of device memory for one prefill round failed every request
of that round; it now returns the sequences to its queue and narrows admission instead. An exam in
which any trial was ended by the server, the environment or the harness decides nothing.

\subsection{Results}

\paragraph{Generation 0.} The base explored the training tasks of single faults, firmware and
mechanical design, and two candidates were trained on its attempts; the generation was decided on
three workloads with the deviations of \cref{app:deviations}. Against the pooled exams the
self-distilling candidate won more held-out tasks than it lost and replaced the base.\cl{G2701} The GRPO candidate gained on both repository workloads but lost on firmware, so the
rule preferred the self-distilling one even before the firmware loss is weighed.\cl{G2702} Nearly
all of the exam was already solved by the base: the gain is made on the few held-out tasks it
solved only sometimes, and most of what first looked like failures of the base were trials the
serving engine or the per-trial limit ended (\cref{sec:generations-harness}).

\paragraph{Generation 1.} Exploring every training task again teaches little: most single-fault
training tasks were solved by every attempt of generation 0, and a group whose attempts all earn the
same reward has no advantage to learn from. Generation 1 therefore practises a frontier:
the training tasks the base did not always solve and the training tasks with several faults it had
never attempted, with the promoted champion. Its GRPO candidate was promoted on a gain that is
not significant.\cl{G2704}

\paragraph{One measurement of all three.} The generation-0 decision compared the candidates with
a base exam taken on an older serving build, part of which the serving engine had ended. When the
base sat the same exam on generation 1's stack, the generation-0 champion no longer beat
it,\cl{G2703} and after two generations no gain over the base is measured.\cl{G2704} The exams are
near their ceiling for this model: the base already solves about nine held-out tasks in ten,\cl{G2703} so the
few tasks it solves only sometimes carry the whole comparison, and two promotions on small
positive means were not enough to move it. The generations did change the model, in one direction: on the tasks both solved, the
generation-1 champion needed about a tenth fewer completion tokens per solved task than the base on every
workload.\cl{G2705} That was found after the data and over several comparisons, so it is exploratory; it
was registered as a co-primary outcome for later generations of these workloads,\cl{G2706} and no later
run of this paper measured it, because the series that followed moved to a task with headroom for a
small model.
Self-improvement in pass rate across generations is therefore not shown by these runs; what they establish is the loop itself, run end to end without a stronger
model, and the conditions an exam must meet before its decision means anything
(\cref{sec:generations-harness}).
